\documentclass[12pt,a4paper]{article}

% ============================================================
% PAQUETES
% ============================================================
\usepackage[utf8]{inputenc}
\usepackage[T1]{fontenc}
\usepackage[spanish,es-tabla,es-nodecimaldot]{babel}
\usepackage{lmodern}
\usepackage{microtype}

\usepackage[a4paper,margin=2.5cm,top=2.5cm,bottom=2.5cm]{geometry}
\usepackage{graphicx}
\usepackage{float}
\usepackage{amsmath,amssymb}
\usepackage{fancyhdr}
\usepackage{setspace}
\usepackage{booktabs}
\usepackage{array}
\usepackage{tabularx}
\usepackage{longtable}
\usepackage{xcolor}
\usepackage{enumitem}
\usepackage{titlesec}
\usepackage{tcolorbox}
\usepackage{hyperref}
\usepackage{caption}
\usepackage{parskip}
\usepackage{listings}

% ============================================================
% COLORES Y ESTILOS
% ============================================================
\definecolor{azulCobalto}{HTML}{1F3A93}
\definecolor{azulClaro}{HTML}{4A6FA5}
\definecolor{grisSuave}{HTML}{F2F4F8}
\definecolor{grisOscuro}{HTML}{3A3A3A}

\hypersetup{
    colorlinks=true,
    linkcolor=azulCobalto,
    urlcolor=azulCobalto,
    citecolor=azulCobalto,
    pdftitle={Reporte del proyecto - Orbit DevOps Task Simulator},
    pdfauthor={Axel Tadeo Diaz Flores},
    pdfsubject={Sistemas Distribuidos},
    pdfkeywords={DevOps, lenguaje natural, OpenAI, Flask, Next.js, SSE}
}

% Formato de secciones
\titleformat{\section}
    {\color{azulCobalto}\Large\bfseries}
    {\thesection.}{0.6em}{}
\titleformat{\subsection}
    {\color{azulClaro}\large\bfseries}
    {\thesubsection.}{0.5em}{}
\titleformat{\subsubsection}
    {\color{grisOscuro}\normalsize\bfseries}
    {}{0em}{}

\titlespacing*{\section}{0pt}{1.4em}{0.6em}
\titlespacing*{\subsection}{0pt}{1em}{0.4em}

% Encabezado y pie de página
\pagestyle{fancy}
\fancyhf{}
\fancyhead[L]{\textcolor{azulCobalto}{\small\textbf{Sistemas Distribuidos}}}
\fancyhead[R]{\textcolor{grisOscuro}{\small Orbit --- Simulador de tareas DevOps}}
\fancyfoot[C]{\thepage}
\renewcommand{\headrulewidth}{0.4pt}
\renewcommand{\headrule}{\hbox to\headwidth{\color{azulClaro}\leaders\hrule height \headrulewidth\hfill}}
\renewcommand{\footrulewidth}{0pt}

% Configuración de listas
\setlist[itemize]{leftmargin=1.5em,itemsep=2pt,topsep=4pt}
\setlist[enumerate]{leftmargin=1.8em,itemsep=2pt,topsep=4pt}

% Espaciado
\setlength{\parskip}{0.6em}
\setlength{\parindent}{0pt}
\onehalfspacing

% Caja informativa
\newtcolorbox{cajainfo}[1][]{
    colback=grisSuave,
    colframe=azulCobalto,
    boxrule=0.6pt,
    arc=2mm,
    left=8pt,right=8pt,top=6pt,bottom=6pt,
    #1
}

% Estilo para bloques de código
\lstset{
    basicstyle=\ttfamily\small,
    backgroundcolor=\color{grisSuave},
    frame=single,
    rulecolor=\color{azulClaro},
    breaklines=true,
    showstringspaces=false,
    keywordstyle=\color{azulCobalto}\bfseries,
    commentstyle=\color{grisOscuro}\itshape,
    stringstyle=\color{azulClaro},
    language=Python,
    numbers=left,
    numberstyle=\tiny\color{grisOscuro},
    xleftmargin=1.5em,
    framexleftmargin=1.5em
}

\begin{document}

% ============================================================
% PORTADA
% ============================================================
\begin{titlepage}
\thispagestyle{empty}

\noindent
\begin{minipage}{0.5\textwidth}
    \IfFileExists{facultad.jpg}{\includegraphics[width=3cm]{facultad.jpg}}{%
        \fbox{\parbox[c][3cm][c]{3cm}{\centering\small Facultad de\\ Ciencias}}}
\end{minipage}%
\begin{minipage}{0.5\textwidth}
    \raggedleft
    \IfFileExists{uabc.png}{\includegraphics[width=3cm]{uabc.png}}{%
        \fbox{\parbox[c][3cm][c]{3cm}{\centering\small UABC}}}
\end{minipage}

\vspace{1cm}

\begin{center}
\textcolor{azulClaro}{\rule{0.85\textwidth}{1.2pt}}

\vspace{0.3cm}

{\large\textsc{Universidad Autónoma de Baja California}}\\[0.2cm]
{\normalsize Facultad de Ciencias}

\vspace{0.3cm}

\textcolor{azulClaro}{\rule{0.85\textwidth}{0.6pt}}

\vspace{1cm}

{\Large\textbf{Sistemas Distribuidos}}

\vspace{2.5cm}

{\Huge\bfseries\textcolor{azulCobalto}{Simulador de tareas técnicas para gestión de infraestructura Dev-Ops}}\\[0.5cm]

\vspace{3cm}

\begin{tabular}{rl}
\textbf{Integrantes:} & Cristian Adrian Vargas Arredondo \\
                 & Jesus Antonio Sauceda Cazares \\
                 & Francisco Gabriel Parga Torres \\
                 & Axel Tadeo Díaz Flores \\[0.4em]
\textbf{Docentes:}    & Mary Sol Aguirre Arzate \\[0.4em]
\textbf{Fecha:}       & Mayo 2026 \\
\end{tabular}

\vfill

\textcolor{azulClaro}{\rule{0.85\textwidth}{1.2pt}}\\[0.2cm]
{\footnotesize Ensenada, Baja California --- Semestre 2026-1}
\end{center}

\end{titlepage}

% ============================================================
% ÍNDICE
% ============================================================
\tableofcontents
\newpage

% ============================================================
% 1. INTRODUCCIÓN
% ============================================================
\section{Introducción}

El presente reporte documenta el desarrollo de \textbf{Orbit}, un simulador de
tareas técnicas para la gestión de infraestructura DevOps. El proyecto consiste
en un \textit{middleware} potenciado por inteligencia artificial cuyo objetivo
es traducir instrucciones expresadas en lenguaje natural a artefactos
ejecutables concretos: comandos o \textit{scripts} de shell (Bash/POSIX) y
consultas SQL.

La motivación del sistema parte de un problema recurrente en la administración
de sistemas y la automatización: existe una brecha entre la \textit{intención}
de alto nivel del operador (``muéstrame los procesos que más memoria consumen'')
y el \textit{comando} de bajo nivel que la materializa
(\texttt{ps aux | sort -k4 -rn | head}). Orbit busca cerrar esa brecha,
permitiendo que un usuario describa lo que desea lograr y reciba a cambio un
artefacto listo para revisar y ejecutar, junto con una explicación técnica
breve.

A diferencia de un simple generador de texto, Orbit implementa un flujo
completo de tres fases —\textbf{Generar}, \textbf{Revisar/Refinar} y
\textbf{Ejecutar}— con \textit{streaming} de resultados en tiempo real y un
conjunto de salvaguardas de seguridad que evitan la ejecución de comandos
potencialmente destructivos.

\begin{cajainfo}
\textbf{En una frase:} Orbit convierte lenguaje natural en \textit{scripts}
DevOps ejecutables, los somete a revisión humana y validación de seguridad, y
los ejecuta en un \textit{subproceso} controlado mostrando la salida paso a
paso.
\end{cajainfo}

% ============================================================
% 2. OBJETIVOS
% ============================================================
\section{Objetivos del proyecto}

\subsection{Objetivo general}

Diseñar e implementar un sistema cliente-servidor que, mediante un modelo de
lenguaje, traduzca instrucciones en lenguaje natural a comandos o consultas
ejecutables, ofreciendo un flujo seguro de generación, refinamiento y ejecución
con retroalimentación visual en tiempo real.

\subsection{Objetivos específicos}

\begin{itemize}
    \item Construir una API REST en Flask que exponga las operaciones de
    generación, refinamiento y ejecución de \textit{scripts}.
    \item Integrar un modelo de lenguaje de OpenAI mediante \textit{prompts} de
    sistema estrictos que garanticen una salida estructurada en formato JSON.
    \item Implementar un motor de ejecución local con validaciones de seguridad
    y verificación previa de la disponibilidad de las herramientas requeridas.
    \item Transmitir el progreso de la ejecución en tiempo real al cliente
    mediante \textit{Server-Sent Events} (SSE).
    \item Desarrollar una interfaz web moderna en Next.js con estética de
    terminal, animaciones fluidas y una máquina de estados que modele el flujo
    de la aplicación.
\end{itemize}

% ============================================================
% 3. DESCRIPCIÓN GENERAL
% ============================================================
\section{Descripción general del sistema}

Orbit se compone de dos aplicaciones independientes que se comunican mediante
HTTP:

\begin{itemize}
    \item Un \textbf{backend} en Python/Flask que actúa como cerebro del
    sistema: recibe las peticiones, dialoga con el modelo de lenguaje, valida la
    seguridad de los \textit{scripts} y los ejecuta en la máquina anfitriona.
    \item Un \textbf{frontend} en Next.js que provee la experiencia de usuario:
    captura la instrucción, presenta el \textit{script} propuesto, permite
    refinarlo o aprobarlo, y renderiza la ejecución y sus resultados.
\end{itemize}

El sistema está diseñado para correr de forma local sobre macOS, empleando un
shell Bash con utilidades POSIX/BSD. Las consultas SQL se ejecutan de manera
simulada (\textit{mock}), pues no hay una base de datos conectada; esto permite
completar el flujo de la interfaz sin efectos colaterales.

% ============================================================
% 4. STACK TECNOLÓGICO
% ============================================================
\section{Stack tecnológico}

A continuación se detallan las tecnologías empleadas en cada capa del sistema,
con las versiones efectivamente utilizadas en el proyecto.

\subsection{Backend}

\begin{itemize}
    \item \textbf{Lenguaje:} Python 3.x.
    \item \textbf{Framework web:} Flask, para exponer la API REST y el
    \textit{streaming} mediante SSE.
    \item \textbf{CORS:} \texttt{flask-cors}, configurado para aceptar
    peticiones desde \texttt{http://localhost:3000}.
    \item \textbf{Variables de entorno:} \texttt{python-dotenv}, que carga el
    archivo \texttt{.env} ubicado en la raíz del proyecto.
    \item \textbf{Integración de IA:} SDK oficial de OpenAI (\texttt{openai}),
    con el modelo \texttt{gpt-4o-mini} por defecto.
    \item \textbf{Ejecución:} módulos estándar \texttt{subprocess},
    \texttt{shutil} y \texttt{re} para correr, validar e inspeccionar los
    \textit{scripts}.
\end{itemize}

\subsection{Frontend}

\begin{itemize}
    \item \textbf{Framework:} Next.js 16 (App Router).
    \item \textbf{Biblioteca de UI:} React 19.
    \item \textbf{Lenguaje:} TypeScript 5.
    \item \textbf{Estilos:} Tailwind CSS v4.
    \item \textbf{Animaciones:} Framer Motion 12.
    \item \textbf{Iconografía:} Lucide React.
    \item \textbf{Composición de clases:} \texttt{clsx} y \texttt{tailwind-merge},
    combinados en la utilidad \texttt{cn(...)}.
\end{itemize}

\subsection{Inteligencia artificial}

\begin{itemize}
    \item \textbf{Proveedor:} OpenAI.
    \item \textbf{Modelo por defecto:} \texttt{gpt-4o-mini}, configurable
    mediante la variable de entorno \texttt{OPENAI\_MODEL}.
    \item \textbf{Formato de respuesta:} modo JSON
    (\texttt{response\_format: json\_object}) con temperatura \texttt{0.2} para
    favorecer salidas deterministas y consistentes.
\end{itemize}

\begin{table}[H]
\centering
\caption{Resumen del stack tecnológico por capa.}
\begin{tabularx}{\textwidth}{@{}l l X@{}}
\toprule
\textbf{Capa} & \textbf{Tecnología} & \textbf{Propósito} \\
\midrule
Backend   & Python 3 / Flask        & API REST y motor de ejecución \\
Backend   & flask-cors              & Política de origen cruzado \\
Backend   & python-dotenv           & Carga de variables de entorno \\
Backend   & openai                  & Cliente del modelo de lenguaje \\
IA        & gpt-4o-mini             & Traducción de lenguaje natural a código \\
Frontend  & Next.js 16 / React 19   & Aplicación web y enrutamiento \\
Frontend  & TypeScript 5            & Tipado estático \\
Frontend  & Tailwind CSS v4         & Sistema de estilos \\
Frontend  & Framer Motion           & Animaciones e interacciones \\
Frontend  & Lucide React            & Iconos \\
\bottomrule
\end{tabularx}
\end{table}

% ============================================================
% 5. ARQUITECTURA
% ============================================================
\section{Arquitectura y estructura del proyecto}

El repositorio se organiza en dos directorios principales, además del archivo de
configuración de entorno en la raíz.

\subsection{Estructura de directorios}

\begin{itemize}
    \item \texttt{/backend}: lógica de la aplicación Flask.
    \begin{itemize}
        \item \texttt{app.py}: punto de entrada, rutas de la API, motor de
        ejecución y verificaciones de seguridad. Corre en el puerto
        \texttt{8000}.
        \item \texttt{requirements.txt}: dependencias de Python
        (\texttt{flask}, \texttt{flask-cors}, \texttt{python-dotenv},
        \texttt{openai}).
    \end{itemize}
    \item \texttt{/frontend}: aplicación Next.js.
    \begin{itemize}
        \item \texttt{/app}: páginas y \textit{layouts} (App Router). La página
        principal \texttt{page.tsx} contiene una máquina de estados basada en
        fases.
        \item \texttt{/components}: componentes de interfaz modulares.
        \item \texttt{/lib}: utilidades compartidas —\texttt{cn.ts} (fusión de
        clases) y \texttt{api.ts} (cliente HTTP tipado con analizador de SSE).
    \end{itemize}
    \item \texttt{/.env}: ubicado en la raíz. Contiene \texttt{OPENAI\_API\_KEY}
    (obligatoria) y, opcionalmente, \texttt{OPENAI\_MODEL} (por defecto
    \texttt{gpt-4o-mini}) y \texttt{EXECUTION\_TIMEOUT} (por defecto
    \texttt{20} segundos).
\end{itemize}

\subsection{Componentes del frontend}

\begin{itemize}
    \item \textbf{PromptInput}: área de texto con ejemplos sugeridos y envío
    mediante el atajo de teclado Cmd+Enter.
    \item \textbf{ScriptPreview}: muestra el \textit{script} generado junto con
    las acciones Ejecutar, Ajustar y Descartar.
    \item \textbf{TerminalBlock}: bloque de código con numeración de líneas y
    resaltado ligero por expresiones regulares para \texttt{bash}, \texttt{sh}
    y \texttt{sql}.
    \item \textbf{ExecutionStream}: progreso de la ejecución paso a paso en
    tiempo real, alimentado por SSE.
    \item \textbf{ResultsPanel}: presenta el resultado con \texttt{stdout} y
    \texttt{stderr}, y un conmutador entre vista formateada y vista cruda.
    \item \textbf{OutputRenderer}: formateo inteligente de la salida (detecta la
    salida de \texttt{ps aux} y la convierte en una tabla de procesos
    interactiva; en otro caso, vista de líneas numeradas).
    \item \textbf{GeneratingSkeleton}: marcador de carga con efecto
    \textit{shimmer} durante la generación.
    \item \textbf{ErrorBanner}: notificación de error descartable.
    \item \textbf{Sidebar}: barra lateral con el historial de sesiones, puntos
    de estado y marcas de tiempo relativas.
\end{itemize}

% ============================================================
% 6. API
% ============================================================
\section{Interfaz de programación (API)}

El backend expone cuatro \textit{endpoints} bajo el prefijo \texttt{/api}.

\begin{table}[H]
\centering
\caption{Endpoints de la API REST.}
\begin{tabularx}{\textwidth}{@{}l l X@{}}
\toprule
\textbf{Método} & \textbf{Ruta} & \textbf{Descripción} \\
\midrule
\texttt{POST} & \texttt{/api/generate} & Genera un \textit{script} a partir de un \textit{prompt}. \\
\texttt{POST} & \texttt{/api/refine}   & Revisa un \textit{script} previo según el \textit{feedback} del usuario. \\
\texttt{POST} & \texttt{/api/execute}  & Ejecuta el \textit{script} y transmite eventos SSE. \\
\texttt{GET}  & \texttt{/api/health}   & Verifica el estado del servicio y el modelo activo. \\
\bottomrule
\end{tabularx}
\end{table}

\begin{itemize}
    \item \texttt{/api/generate}: recibe \texttt{\{ "prompt": "..." \}} y
    devuelve \texttt{\{ "id", "script", "explanation", "language" \}}.
    \item \texttt{/api/refine}: recibe el \textit{prompt} original, el
    \textit{script} previo y la retroalimentación, y devuelve un \textit{script}
    revisado completo (no un \textit{diff}).
    \item \texttt{/api/execute}: recibe \texttt{\{ "script", "language" \}} y
    devuelve un flujo SSE con eventos \texttt{start}, \texttt{step} (uno por
    cada fase) y \texttt{done}.
    \item \texttt{/api/health}: devuelve \texttt{\{ "status": "ok", "model":
    "..." \}}.
\end{itemize}

% ============================================================
% 7. LÓGICA DE IA
% ============================================================
\section{Lógica de inteligencia artificial}

El núcleo inteligente del sistema reside en el diseño de los \textit{prompts} de
sistema y en el procesamiento de la respuesta del modelo.

\subsection{Prompt de sistema}

El \textit{prompt} de sistema, redactado en español, instruye al modelo a actuar
como un experto en automatización DevOps. Sus reglas más relevantes son:

\begin{itemize}
    \item Devolver siempre un \textit{script} ejecutable, válido y no vacío.
    \item Generar sintaxis BSD/POSIX compatible con macOS, evitando extensiones
    propias de GNU/Linux (por ejemplo, \texttt{ps aux -{}-sort} o
    \texttt{grep -P}).
    \item Preferir herramientas POSIX estándar y comandos seguros e
    idempotentes, evitando operaciones destructivas salvo petición explícita.
    \item Generar comandos de herramientas especializadas (Docker, kubectl,
    systemctl, Terraform, etc.) únicamente cuando el usuario las menciona de
    forma explícita.
    \item Responder con un objeto JSON que contenga exactamente las claves
    \texttt{script}, \texttt{explanation} (una oración técnica en español) y
    \texttt{language} (uno de \texttt{bash}, \texttt{sh} o \texttt{sql}), sin
    texto conversacional.
\end{itemize}

\subsection{Refinamiento}

El \textit{endpoint} \texttt{/api/refine} emplea un \textit{prompt} de sistema
distinto que recibe la solicitud original, el \textit{script} anterior y el
\textit{feedback} del usuario, y produce una versión revisada completa
manteniendo el mismo contrato JSON.

\subsection{Análisis tolerante de la salida}

La función \texttt{\_parse\_model\_output} interpreta la respuesta del modelo de
forma robusta: primero intenta analizar el JSON; si falla, busca un bloque de
código delimitado por comillas invertidas; y, como último recurso, trata todo el
texto como el \textit{script}. Esto evita que pequeñas desviaciones en el
formato rompan el flujo.

% ============================================================
% 8. SEGURIDAD Y EJECUCIÓN
% ============================================================
\section{Seguridad y motor de ejecución}

Dado que el sistema ejecuta código en la máquina anfitriona, la seguridad es una
preocupación central. El motor de ejecución aplica varias salvaguardas antes y
durante la ejecución.

\subsection{Validación de patrones peligrosos}

Antes de ejecutar, el \textit{script} se contrasta contra una lista de patrones
peligrosos mediante expresiones regulares. Entre ellos se incluyen
\texttt{rm -rf /}, las \textit{fork bombs} clásicas, \texttt{mkfs}, escrituras
crudas con \texttt{dd} hacia dispositivos, \texttt{shutdown}, \texttt{reboot} y
redirecciones hacia \texttt{/dev/sda}. Si se detecta una coincidencia, la
ejecución se rechaza de inmediato.

\subsection{Verificación de herramientas}

A continuación, el sistema analiza el \textit{script} para detectar los comandos
externos referenciados y comprueba —mediante \texttt{shutil.which}— que estén
instalados. Si falta alguna herramienta (por ejemplo, \texttt{docker} o
\texttt{kubectl}), se informa al usuario antes de gastar un \textit{subproceso},
sugiriendo reformular la tarea con utilidades disponibles.

\subsection{Ejecución controlada}

Los \textit{scripts} de shell se ejecutan mediante
\texttt{subprocess.run(["/bin/bash", "-c", script])} con un \textit{timeout}
configurable (20 segundos por defecto). Las consultas SQL se ejecutan de forma
simulada, devolviendo una salida controlada. Cada fase del proceso —validación,
verificación, preparación del entorno, ejecución y recolección de salida— se
emite como un evento SSE, de modo que el cliente puede mostrar el avance en
tiempo real.

% ============================================================
% 9. FLUJO DE TRABAJO
% ============================================================
\section{Flujo de trabajo}

El ciclo completo de interacción, desde la instrucción hasta el resultado, sigue
estos pasos:

\begin{enumerate}
    \item El usuario introduce una tarea en lenguaje natural en
    \textbf{PromptInput}.
    \item El frontend invoca \texttt{POST /api/generate} con el \textit{prompt}.
    \item El backend consulta a OpenAI con el \textit{prompt} de sistema y el
    modo de respuesta JSON.
    \item OpenAI devuelve \texttt{\{ "script", "explanation", "language" \}}.
    \item El backend valida o analiza la salida y la regresa con un \texttt{id}
    generado.
    \item El frontend entra en la fase \texttt{suggested}: muestra el
    \textit{script} en un \textbf{TerminalBlock}, la explicación y tres
    acciones:
    \begin{itemize}
        \item \textbf{Ejecutar}: aprueba y ejecuta el \textit{script}.
        \item \textbf{Ajustar}: abre un área de refinamiento que envía el
        \textit{feedback} a \texttt{/api/refine} y reemplaza el \textit{script}.
        \item \textbf{Descartar}: regresa al estado inicial.
    \end{itemize}
    \item Al aprobar, el frontend llama a \texttt{POST /api/execute} y entra en
    la fase \texttt{executing}.
    \item El backend transmite eventos SSE: valida seguridad, comprueba
    herramientas, ejecuta el \textit{script} y recolecta la salida.
    \item El frontend renderiza cada paso en \textbf{ExecutionStream} en tiempo
    real.
    \item Al finalizar, \textbf{ResultsPanel} muestra el código de salida, la
    duración, \texttt{stdout} (formateado o crudo) y \texttt{stderr}.
    \item Las sesiones se registran en la \textbf{Sidebar} con indicadores de
    estado.
\end{enumerate}

\subsection{Máquina de estados}

La página principal del frontend modela el flujo mediante una unión
discriminada de fases:
\[
\texttt{idle} \rightarrow \texttt{generating} \rightarrow \texttt{suggested}
\rightarrow \texttt{executing} \rightarrow \texttt{finished}
\]
La fase \texttt{suggested} admite además un subestado opcional de refinamiento
(\texttt{refining}).

% ============================================================
% 10. DISEÑO DE LA INTERFAZ
% ============================================================
\section{Diseño de la interfaz}

La interfaz adopta una estética moderna de terminal/DevOps:

\begin{itemize}
    \item \textbf{Modo oscuro} exclusivo, con fondo \texttt{\#09090b} y fuentes
    monoespaciadas para el código.
    \item \textbf{Glassmorphism}: tarjetas con \texttt{backdrop-blur-sm},
    bordes sutiles (\texttt{border-white/10}) y color de acento esmeralda.
    \item \textbf{Animaciones} con Framer Motion: transiciones con
    \texttt{AnimatePresence}, entradas escalonadas, deslizamiento y aparición
    gradual, y estados de carga.
    \item \textbf{Fondo} construido con pseudo-elementos \texttt{::before}
    (gradientes radiales en índigo y esmeralda) y \texttt{::after} (rejilla CSS
    de 48 px).
    \item \textbf{Texto} de la interfaz íntegramente en español.
\end{itemize}

% ============================================================
% 11. EJECUCIÓN LOCAL
% ============================================================
\section{Ejecución local}

Para levantar el sistema en un entorno de desarrollo se requieren dos terminales.

\subsection{Backend (puerto 8000)}

\begin{lstlisting}[language=bash]
cd backend
source .venv/bin/activate
python app.py
\end{lstlisting}

\subsection{Frontend (puerto 3000)}

\begin{lstlisting}[language=bash]
cd frontend
npm run dev
\end{lstlisting}

Una vez ambos servicios estén activos, la aplicación es accesible en
\texttt{http://localhost:3000} y consume la API en
\texttt{http://localhost:8000}.

% ============================================================
% 12. CONCLUSIONES
% ============================================================
\section{Conclusiones}

Orbit demuestra cómo un modelo de lenguaje puede integrarse de forma segura y
útil dentro de un flujo de trabajo de administración de sistemas. Más allá de la
mera generación de texto, el proyecto aporta valor mediante tres decisiones de
diseño clave: la \textbf{revisión humana obligatoria} antes de cualquier
ejecución, las \textbf{salvaguardas de seguridad} que bloquean comandos
destructivos y verifican la disponibilidad de herramientas, y la
\textbf{retroalimentación en tiempo real} que transparenta cada fase del
proceso.

Desde el punto de vista de los sistemas distribuidos, el proyecto ejercita la
comunicación cliente-servidor sobre HTTP, el manejo de políticas de origen
cruzado (CORS), el \textit{streaming} unidireccional mediante \textit{Server-Sent
Events} y la coordinación de una máquina de estados entre dos procesos
independientes.

Como trabajo futuro, el sistema podría extenderse con la conexión a una base de
datos real para las consultas SQL, la persistencia del historial de sesiones, la
autenticación de usuarios y la ejecución de los \textit{scripts} dentro de
contenedores aislados para reforzar aún más las garantías de seguridad.

\end{document}
